1. Metabolic Modeling & DEB (Direct Foundation)

These are essential to position your work as an extension of existing theoretical frameworks.

    Eriksen, N. T. (2022). Dynamic modelling of feed assimilation, growth, lipid accumulation, and CO2 production in black soldier fly larvae.

        Relevance: This is the core reference you already cite. It establishes the theoretical link between respiration and growth.

    Eriksen, N. T. (2024). Metabolic performance and feed efficiency of black soldier fly larvae.

        Relevance: A recent update that further explores metabolic flows, reinforcing the biological validity of using CO2 as a proxy for efficiency.

    Hansen, R. J., et al. (2023). Metabolic performance of black soldier fly larvae during growth on different diets.

        Relevance: Validates that metabolic rates vary significantly with diet quality, supporting your D1 vs. D4 comparison.

2. Real-Time Monitoring via Respirometry (The "Virtual Sensor" Concept)

This is the most critical category. These papers prove that your idea of inferring biomass from gas is a current "hot topic."

    [Authors TBA] (2025 Preprint/Article). Using gas exchange measurements to monitor growth, energy expenditure, and body composition of black soldier fly larvae.

        Relevance: Extremely high. This study explicitly proposes "integrating respirometry into rearing systems to enable real-time monitoring of insect biomass yield." It confirms that your "Virtual Sensor" is a scientifically valid and novel pursuit.

    Muurmann, A. T., et al. (2025). Growth and metabolic performance of house fly and black soldier fly.

        Relevance: Compares metabolic parameters and substrate conversion, providing a broader context for your physiological assumptions.

3. Greenhouse Gas Emissions (GWP Context)

Use these to benchmark your results regarding the "environmental cost" of the slow diet (D4).

    Mertenat, A., et al. (2019). Bioconversion efficiencies, greenhouse gas and ammonia emissions during black soldier fly rearing—A mass balance approach.

        Relevance: The seminal paper on BSF emissions. It provides the baseline values for CO2, CH4, and N2O that you should compare your integral results against.

    [Authors TBA] (2024). Black soldier fly larvae mitigate greenhouse gas emissions in the treatment of organic waste.

        Relevance: Explains the microbial mechanism (e.g., methanogens) behind the emissions. This supports your discussion on how larvae reduce the "anaerobic window" (your D1 argument).

4. PINNs in Bioprocesses (Methodological Validation)

Since there are few PINN papers specifically on BSF, cite these to justify the method in broader biotechnology.

    Cui, et al. (2024/2025). Physics-informed neural networks for multi-stage Koopman modeling of microbial fermentation.

        Relevance: Demonstrates the use of PINNs to model fermentation dynamics (like your CH4 production) where data is sparse and physical laws (mass balance) must be respected.

    [Recent Research] (2025). Bioprocess model-predictive control with physics-informed neural networks.

        Relevance: Shows the trend of moving from "passive modeling" to "active control" (monitoring) using PINNs, which aligns with your thesis goal.
